%\vspace{-1em}
\section{Conclusion and Limitations}
%\vspace{-1em}
By formulating manipulation as mode-conditioned dynamics, we illustrate why Contrastive Reinforcement Learning will struggle in manipulation. \methodname{} offers a practical, simple solution that provides an inductive bias towards this complexity, resulting in improved performance across interaction-rich tasks---especially those with rarer or more complex interactions, such as hard dynamic ball interactions, pick-and-place, and real-world air hockey. As robot learning moves towards more complex settings, handling this complexity will become increasingly crucial.

\textbf{Limitations:} Despite the promising success of \methodname{}, it remains an incomplete initial step toward leveraging the insight of mode-conditioned dynamics in robot learning. First, \methodshort{} still represents multimodal manipulation dynamics with a unimodal model---it only changes sampling. Second, CRL is one form of unsupervised RL, with limited representational power, but \methodshort{} can be applied to other unsupervised RL methods. Third, \methodshort{} assumes a factorized state and heuristic interaction indicators---integrating \methodshort{} with an end-to-end perception system is key to future systems. Fourth, CRL incurs a high sample-efficiency cost, preventing it from real-world application---our failures stem from the sim-to-real gap. Thus, future directions involve both algorithmic modifications and better robot systems for transferring core knowledge to the agent.